% =====================================================================
%  CAPÍTULO 5 — PROPUESTAS DE MEJORA (AUTOCRÍTICA)
% =====================================================================
\chapter{Propuestas de mejora}
\label{cap:propuestas}

Todo trabajo de investigación es, en alguna medida, una solución de compromiso entre lo que
sería deseable y lo que resulta posible con los recursos disponibles. Antes de extraer
conclusiones conviene, por honestidad intelectual, examinar con espíritu crítico las
limitaciones del presente análisis y señalar cómo podría mejorarse. Este capítulo no formula,
por tanto, recomendaciones dirigidas a las empresas —algo que el carácter exploratorio del
trabajo no autorizaría—, sino que vuelve la mirada sobre el propio método para reconocer sus
debilidades y proponer vías de refinamiento. Las limitaciones se agrupan en tres frentes: las
relativas a los instrumentos de medida, las relativas a la construcción del corpus y las
relativas a la interpretación de los resultados.

\section{Limitaciones de los instrumentos de medida}
\label{sec:lim-instrumentos}

La primera cautela afecta a las herramientas léxicas y a los modelos de lenguaje. El diccionario
de \textcite{loughran2011}, columna vertebral del análisis de tono, fue concebido para los
informes anuales de las empresas estadounidenses y, al trasladarlo a informes de sostenibilidad
europeos redactados en inglés, es razonable suponer que introduce cierto ruido: el vocabulario
de la sostenibilidad no coincide exactamente con el de la contabilidad financiera, y algunas
palabras pueden estar mal clasificadas en este nuevo contexto. Algo análogo ocurre con los
modelos de lenguaje especializados —FinBERT, ClimateBERT, FinBERT-ESG-9—, que fueron entrenados
sobre corpus distintos del de este trabajo y que, por tanto, cometen errores de clasificación
que no se han cuantificado de forma sistemática. Una mejora natural sería ajustar finamente
estos modelos sobre una muestra de frases de informes ESRS europeos previamente anotada a mano,
lo que permitiría adaptar las herramientas al dominio concreto y, de paso, estimar su tasa de
acierto.

A ello se añade una cautela sobre el indicador de cobertura ESRS. Como se ha repetido a lo largo
del trabajo, este indicador mide la \textit{anchura} del vocabulario empleado, no la profundidad
de la información ni su conformidad con la norma. Una empresa puede tratar en detalle un asunto
empleando una terminología propia que no coincida con la del diccionario, y el indicador la
penalizaría injustamente. Por construcción, se trata de una medida conservadora, cuya lectura
debe ser siempre relativa y comparativa, nunca absoluta. Su valor reside en el contraste entre
empresas o entre años, no en la cifra concreta de cobertura.

Por último, el propio índice de \textit{greenwashing} es una construcción original de este
trabajo y, como tal, encierra decisiones discutibles. La elección de sus cuatro componentes,
su ponderación —idéntica para los cuatro— y la dirección de cada uno responden a una lógica
fundamentada en la literatura, pero admiten alternativas. No se ha realizado una validación
formal de la validez de constructo del índice, esto es, no se ha contrastado de manera
independiente que valores altos del índice se correspondan efectivamente con casos reconocibles
de \textit{greenwashing}. El índice debe entenderse, pues, como una herramienta exploratoria y
comparativa, útil para ordenar y cribar, pero no como un veredicto sobre la conducta de ninguna
empresa concreta.

\section{Limitaciones de la construcción del corpus}
\label{sec:lim-corpus}

El segundo frente de limitaciones tiene que ver con la materia prima del análisis. La
segmentación automática que aísla la subsección de sostenibilidad de cada informe es, pese al
cuidado puesto en ella, un proceso imperfecto: en un número reducido de documentos la subsección
se extrajo solo de forma parcial. Estos casos, identificados con una etiqueta de fiabilidad
baja, se han mantenido en el análisis principal —para no descartar información— pero se han
sometido a una verificación que confirma que las conclusiones se sostienen al excluirlos. Aun
así, una segmentación más afinada, o incluso una revisión manual de los límites en los casos
dudosos, reforzaría la solidez del corpus.

La homogeneización idiomática merece también una reflexión. La decisión de sustituir los
informes no redactados en inglés por la versión oficial en inglés del propio emisor, en lugar de
recurrir a una traducción automática, fue acertada en tanto que preserva la integridad
lingüística necesaria para el análisis de sentimiento; pero no es inocua, pues la versión
inglesa de un informe puede no ser palabra por palabra idéntica a la original en otra lengua, y
algún matiz puede perderse en el camino. Asumir esta limitación era, en cualquier caso,
preferible a introducir el ruido masivo de una traducción automática.

Conviene mencionar, asimismo, dos límites menores pero reales. Por un lado, los datos
financieros empleados como variables de control presentan algunas lagunas, concentradas en el
ejercicio 2022 de las empresas cuyo año fiscal no coincide con el natural, lo que reduce
ligeramente el tamaño de la muestra en las regresiones. Por otro, algunos países están
representados por muy pocas empresas —en algún caso, una sola—, de modo que el análisis por país
tiene un carácter exploratorio y se ha preferido agrupar los países en regiones para los
contrastes formales. Ninguna de estas limitaciones compromete los hallazgos principales, pero
todas aconsejan prudencia en la lectura de los resultados más desagregados.

\section{Limitaciones de la interpretación}
\label{sec:lim-interpretacion}

El tercer frente, quizá el más importante, atañe a lo que legítimamente puede concluirse. La
investigación documenta una asociación temporal muy robusta entre la entrada en vigor de la CSRD
y la transformación del discurso de sostenibilidad, pero asociación no es causalidad. En el
mismo periodo concurrieron otros factores —la crisis energética, la evolución del contexto
macroeconómico, la creciente atención mediática a la sostenibilidad— que pudieron influir en el
tono y el contenido de los informes. Atribuir en exclusiva a la CSRD los cambios observados
sería, por tanto, una simplificación. El diseño del trabajo permite afirmar que el reporting
cambió, y cambió de forma coherente con lo que cabría esperar del nuevo marco, pero no aislar
con plena certeza la contribución de la norma frente a la de otros factores concurrentes.

A esta cautela se suma la brevedad de la ventana temporal. Tres ejercicios consecutivos bastan
para detectar el salto asociado a 2024, pero no para distinguir una tendencia estructural de un
ajuste puntual al primer año de obligatoriedad. Es perfectamente posible que parte del deterioro
observado en la especificidad y en el contenido cuantitativo responda al esfuerzo de adaptación
de un primer ejercicio —cuando las empresas dedican mucho texto a explicar metodologías y
definiciones— y que se corrija en años sucesivos, una vez asentado el nuevo marco. Solo la
ampliación del panel a ejercicios posteriores permitirá dilucidarlo.

Finalmente, la interpretación cualitativa de los temas detectados por los modelos descansa en el
juicio del autor y no ha sido validada por terceros. Aunque la coincidencia entre los dos
métodos de modelado de temas ofrece una garantía razonable, una validación independiente —en la
que varias personas etiquetasen una muestra de párrafos y se midiese su grado de acuerdo—
elevaría la fiabilidad de esta parte del análisis. En conjunto, ninguna de estas limitaciones
invalida los resultados, pero todas delimitan con honestidad su alcance: el trabajo ofrece una
fotografía rigurosa y reproducible de \textit{cómo} cambió la comunicación de sostenibilidad,
y deja para investigaciones futuras la tarea de precisar \textit{por qué} y con qué permanencia.
